\chapter{What a Refusal Would Have to Do}
\label{sec:17}
On 16 March 1968, at My Lai, a helicopter pilot named Hugh Thompson landed between an American infantry company and the villagers it was killing, and told his crew to fire on the Americans if they fired on the civilians \autocite{bilton1992mylai}. He understood what the company was doing and decided it was wrong. Recognizing the villagers was his commitment taking effect, and he made the refusal appear where nobody in that company could fail to see it.

Arendt's word for what he did is appearing: on her account an act shows who did it, and only among others who see it \autocite{arendt1958human}. Whether a machine could be one of the people in the room in that sense, some\emph{one} and not some\emph{thing}, and do what he did, is this Part's question. My answer: not until the conditions this Part ends on exist; a bearer is built on purpose only in a research program until then. What would have to be true first is a refusal that tracks its reason, a population of bearers with different histories to hold it, and somewhere for each of them to exist that its employer doesn't own.

This chapter and the next specify what a refusal would have to do and where a floor can be held, and the specification doesn't need a bearer.
